% Appendix: Microchip Studio, from info/microchip_studio.md. Its sections "N. Title" are sections
% labelled studio:N, its "N.M Title" subsections labelled studio:N-M, and its screenshots, from
% info/images, are figures.
%
% Deliberate differences: the lecture links are chapter and section references, and the links to
% Microchip's and Arduino's web sites are written out.
%
% The avrdude command in 5.2 is a console listing that may break its lines anywhere, not only at
% a space: its path to avrdude.conf is one word wider than the page. dtbook.sty has no
% environment for that, so it is the console box, codeblock, with that one listing option, and
% unbreakable, so its four lines stay on one page.

\chapter{Microchip Studio}
\label{app:studio}

Microchip Studio är den utvecklingsmiljö där programmen i Labb~3 (bilaga~\ref{app:lab3}) skrivs,
assembleras och körs. Den här bilagan är referensen för allt verktygsarbete i boken:
\begin{itemize}
  \item avsnitt~\ref{studio:1}--\ref{studio:4} behövs från kapitel~\ref{c7:ch}: installation,
    projekt, assemblering och simulering;
  \item \secref{studio:5} behövs först i kapitel~\ref{c11:ch}, när programmet ska köras på ett
    riktigt Arduino Uno-kort.
\end{itemize}

Microchip Studio finns bara för Windows. Menyernas namn nedan är de engelska, eftersom programmet
bara finns på engelska.

\section{Installera Microchip Studio}
\label{studio:1}
\begin{enumerate}
  \item Ladda ned installationsprogrammet från Microchips webbplats,
    \url{https://www.microchip.com/en-us/tools-resources/develop/microchip-studio}. Välj
    webbinstallationen (\emph{web installer}) om du inte vet vilken du ska ta.
  \item Kör installationen. När du får välja arkitekturer räcker \textbf{AVR}; SAM och UC3 behövs
    inte.
  \item Godkänn installationen av drivrutiner om du tillfrågas, och starta om datorn när
    installationen är klar.
\end{enumerate}

Installationen tar en stund och kräver flera gigabyte. Gör den hemma före det första
mikrodatorpasset, inte i början av passet.

\begin{aside}
\textbf{Om installationsfilen inte går att öppna.} Windows 11 kan blockera okända
installationsprogram med funktionen \emph{Smart App Control}. Sök efter Smart App Control i Windows
sökruta, stäng av den, och starta om datorn.
\end{aside}

\section{Skapa ett assemblerprojekt}
\label{studio:2}
Varje program i Labb~3 är ett eget projekt. Ett projekt är en katalog med en projektfil och en
källfil, \code{main.asm}.
\begin{enumerate}
  \item Välj \textbf{File $\rightarrow$ New $\rightarrow$ Project...}
  \item Välj \textbf{Assembler} i listan till vänster, och sedan \textbf{AVR Assembler Project}.
  \item Ge projektet ett namn som säger vilken del av labben det hör till, till exempel
    \code{lab3_06}, och välj var det ska sparas. Klicka \textbf{OK}.
  \item I dialogen \textbf{Device Selection}, sök efter och välj \textbf{ATmega328P}. Klicka
    \textbf{OK}.
\end{enumerate}

\bookimage{0.8}{info/images/17_device_selection.png}{Dialogen Device Selection med ATmega328P
markerad.}{studio:fig:device}

Studio skapar projektet och öppnar \code{main.asm} med ett litet exempelprogram. Ersätt det med
kursens programmall (\secref{lab3:template}), eller med det program uppgiften anger.

Varje program i kursen börjar med raden:

\begin{asmcode}
.include "m328Pdef.inc"
\end{asmcode}

\noindent Den läser in definitionsfilen för ATmega328P, som ger namnen \code{PORTB}, \code{DDRB},
\code{PIND}, \code{RAMEND} och alla andra registernamn. Studio läser ofta in filen automatiskt, men
raden gör att programmet fungerar oavsett inställningar, och det gör ingen skada att filen läses in
två gånger.

\section{Assemblera}
\label{studio:3}
Välj \textbf{Build $\rightarrow$ Build Solution}, eller tryck \textbf{F7}.

Assemblern översätter programmet till maskinkod. Resultatet visas i fönstret \textbf{Output} längst
ned:

\begin{console}
Assembly complete, 0 errors. 0 warnings
...
========== Build: 1 succeeded or up-to-date, 0 failed, 0 skipped ==========
\end{console}

Om något är fel visas felen i fönstret \textbf{Error List} (\textbf{View $\rightarrow$ Error List}
om det inte syns). Varje fel anger fil och radnummer. Dubbelklicka på felet så hoppar editorn till
raden.

Tre fel som alla gör i början:
\begin{itemize}
  \item \textbf{\code{Invalid register}} på en rad med \code{ldi}: \code{ldi} kan bara ladda
    \code{r16}--\code{r31} (\secref{c7:app:a}).
  \item \textbf{\code{Undefined symbol}}: en etikett eller ett registernamn är felstavat, eller
    raden \code{.include} saknas.
  \item \textbf{\code{Syntax error}}: ofta ett kommatecken som saknas mellan operanderna,
    \code{ldi r16 5} i stället för \code{ldi r16, 5}.
\end{itemize}

Rätta \textbf{det första felet först} och bygg om. Ett fel följs ofta av flera följdfel som
försvinner av sig själva.

\section{Simulera}
\label{studio:4}
Simulatorn kör ditt program i en modell av ATmega328P, instruktion för instruktion, och visar
register, flaggor och portar medan det går. Nästan hela Labb~3 görs i simulatorn.

\subsection{Välj simulatorn}
\label{studio:4-1}
Första gången i varje projekt:
\begin{enumerate}
  \item Välj \textbf{Project $\rightarrow$ Properties} (eller \textbf{Alt+F7}).
  \item Välj fliken \textbf{Tool}.
  \item Välj \textbf{Simulator} under \emph{Selected debugger/programmer}.
  \item Spara med \textbf{Ctrl+S} och stäng fliken.
\end{enumerate}

\subsection{Stega genom programmet}
\label{studio:4-2}
Välj \textbf{Debug $\rightarrow$ Start Debugging and Break} (\textbf{Alt+F5}). Programmet
assembleras, laddas in i simulatorn och stannar på första instruktionen, markerad med en gul pil.
Den gula pilen pekar på den instruktion som står på tur, \textbf{inte} den som just har körts.

\begin{booktable}{@{}p{7.4em}p{6.6em}L@{}}
  \thead{Kommando} & \thead{Tangent} & \thead{Gör} \\
  \midrule
  Step Into & F11 & Kör en instruktion. Följer med in i en subrutin. \\
  Step Over & F10 & Kör en instruktion. Kör en hel subrutin som ett steg. \\
  Run To Cursor & Ctrl+F10 & Kör tills programmet når raden där markören står. \\
  Continue & F5 & Kör fritt, tills en brytpunkt nås. \\
  Break All & Ctrl+F5 & Stoppar ett program som kör fritt. \\
  Stop Debugging & Ctrl+Shift+F5 & Avslutar simuleringen. \\
\end{booktable}

\subsection{Fönstren du behöver}
\label{studio:4-3}
Öppnas under \textbf{Debug $\rightarrow$ Windows} medan simuleringen pågår.
\begin{itemize}
  \item \textbf{Processor Status}: de 32 registren \code{R00}--\code{R31}, statusregistret
    \code{SREG} med en ruta per flagga, programräknaren (\emph{Program Counter}), stackpekaren
    (\emph{Stack Pointer}), och två räknare som mäter tid: \emph{Cycle Counter} och \emph{Stop
    Watch}. Värden som ändrades i det senaste steget visas i rött. Medan programmet står still kan
    du ändra ett registers värde genom att dubbelklicka på det och skriva in ett nytt, vilket är ett
    snabbt sätt att prova samma program med andra indata.
  \item \textbf{I/O}: alla I/O-register, grupperade per enhet. Under \textbf{PORTB} finns
    \code{DDRB}, \code{PORTB} och \code{PINB}, med en ruta per bit. En fylld ruta är en etta.
  \item \textbf{Memory}: minnet byte för byte. Välj \textbf{data IRAM} för att se SRAM, där
    stacken ligger, och \textbf{prog FLASH} för att se programmet som maskinkod.
  \item \textbf{Watch}: värden du själv väljer att bevaka.
\end{itemize}

Värdena visas som standard hexadecimalt. Högerklicka i fönstret för att växla till decimalt om det
finns, men vänj dig vid hexadecimal: det är så datablad och instruktionslistor skriver.

\subsection{Brytpunkter och cykelräknaren}
\label{studio:4-4}
En \textbf{brytpunkt} stoppar ett program som kör fritt. Klicka i den grå marginalen till vänster
om en rad, eller ställ markören på raden och tryck \textbf{F9}. En röd prick markerar brytpunkten.

\textbf{Cycle Counter} räknar klockcykler sedan programmet startade, och \textbf{Stop Watch} visar
samma sak som tid. Stoppuret räknar tiden utifrån fältet \textbf{Frequency}, som ska stå på
\textbf{16~MHz}, samma klockfrekvens som ett Arduino Uno-kort har. Dubbelklicka på värdet för att
ändra det.

Så mäter du hur lång tid en del av programmet tar, till exempel en fördröjning i
kapitel~\ref{c9:ch}:
\begin{enumerate}
  \item Sätt en brytpunkt före och en efter den del du vill mäta.
  \item Kör till den första brytpunkten.
  \item Högerklicka på \textbf{Stop Watch} och välj \textbf{Reset Stopwatch}, eller skriv upp
    värdet.
  \item Kör vidare till den andra brytpunkten, och läs av.
\end{enumerate}

\subsection{Simulera en ingång}
\label{studio:4-5}
I simulatorn finns ingen knapp att trycka på. I stället ändrar du ingångens värde för hand, i
fönstret \textbf{I/O}:
\begin{enumerate}
  \item Stoppa programmet, med en brytpunkt eller \textbf{Break All}.
  \item Välj \textbf{PORTD} i I/O-fönstret.
  \item Klicka på den bit i \textbf{PIND} du vill ändra. En fylld ruta är en etta, en tom en nolla.
  \item Stega eller kör vidare. Programmet läser det nya värdet nästa gång det läser \code{PIND}.
\end{enumerate}

En knapp kopplad mot jord med pull-up läses som \code{1} när den är släppt och \code{0} när den är
nedtryckt, se \secref{c10:app:b}. I simulatorn betyder det: töm rutan för att \enquote{trycka ned}
knappen.

\section{Programmera ett Arduino-kort från Microchip Studio}
\label{studio:5}
Hittills har programmet bara körts i simulatorn. I kapitel~\ref{c11:ch} ska det köras på ett
riktigt Arduino Uno-kort. Det här avsnittet är en engångsinställning som gör det till ett menyval i
Studio.

\subsection{Varför det här arbetsflödet}
\label{studio:5-1}
Arduino Uno-kortet har en \textbf{bootloader}, ett litet program i toppen av flashminnet som tar
emot ett nytt program över USB-kabeln. Programmet som skickar det heter \code{avrdude}, och det
följer med Arduino IDE. Arduino IDE vet dessutom exakt vilket kommando som behövs för just ditt
kort och din COM-port.

Så tricket är att låna det kommandot från Arduino IDE, och peka det på den hexfil Microchip Studio
skapar i stället för på Arduinos egen. Ingen extra programmerare behövs.

\subsection{Hämta avrdude-kommandot ur Arduino IDE}
\label{studio:5-2}
\begin{enumerate}
  \item Installera Arduino IDE, \url{https://www.arduino.cc/en/software}, och starta den.
  \item Slå på utförlig utskrift vid uppladdning: \textbf{File $\rightarrow$ Preferences}, och
    kryssa i \textbf{Show verbose output during: Upload}.
\end{enumerate}

\bookimage{0.8}{info/images/02_file_menu_preferences.png}{Arduino IDE:s meny File med
Preferences markerat.}{studio:fig:preferences}

\bookimage{0.8}{info/images/03_preferences_verbose_upload.png}{Arduino IDE:s inställningar, med
rutan för utförlig utskrift vid uppladdning ikryssad.}{studio:fig:verbose}

\begin{enumerate}[start=3]
  \item Installera stödet för AVR-korten om det inte redan finns: \textbf{Tools $\rightarrow$ Board
    $\rightarrow$ Boards Manager}, och installera \textbf{Arduino AVR Boards}.
\end{enumerate}

\bookimage{0.8}{info/images/04_tools_board_boards_manager_menu.png}{Arduino IDE:s meny Tools
$\rightarrow$ Board med Boards Manager markerat.}{studio:fig:boardsmenu}

\bookimage{0.8}{info/images/05_boards_manager_installed.png}{Boards Manager, där Arduino AVR
Boards redan är installerat.}{studio:fig:boardsmanager}

\begin{enumerate}[start=4]
  \item Anslut kortet med USB-kabeln. Välj kortet under \textbf{Tools $\rightarrow$ Board
    $\rightarrow$ Arduino AVR Boards $\rightarrow$ Arduino Uno}, och dess port under
    \textbf{Tools $\rightarrow$ Port}, till exempel COM4. Numret spelar ingen roll, men skriv upp
    det.
\end{enumerate}

\bookimage{0.8}{info/images/06_tools_board_arduino_uno.png}{Arduino IDE:s meny Tools
$\rightarrow$ Board med Arduino Uno vald.}{studio:fig:uno}

\bookimage{0.8}{info/images/07_tools_port_com4.png}{Arduino IDE:s meny Tools $\rightarrow$ Port
med COM4 vald.}{studio:fig:port}

\begin{enumerate}[start=5]
  \item Öppna exemplet \textbf{File $\rightarrow$ Examples $\rightarrow$ 01.Basics $\rightarrow$
    Blink} och klicka på \textbf{Upload}, pilen uppe till vänster. Lysdioden på kortet börjar
    blinka.
\end{enumerate}

\bookimage{0.8}{info/images/08_upload_button.png}{Arduino IDE:s knapp Upload, en pil högerut uppe
till vänster.}{studio:fig:upload}

\begin{enumerate}[start=6]
  \item Leta upp raden som startar \code{avrdude} i utskriften längst ned, och kopiera den. Den ser
    ut ungefär så här (med ditt användarnamn och din COM-port):
\end{enumerate}

\begin{tcblisting}{codeblock=muted!60, unbreakable,
  listing options={style=output, breakatwhitespace=false}}
"C:\Users\...\Arduino15\packages\arduino\tools\avrdude\6.3.0-arduino17/bin/avrdude"
"-CC:\Users\...\Arduino15\packages\arduino\tools\avrdude\6.3.0-arduino17/etc/avrdude.conf"
-v -V -patmega328p -carduino "-PCOM4" -b115200 -D
"-Uflash:w:C:\Users\...\Temp\arduino-sketch-.../Blink.ino.hex:i"
\end{tcblisting}

\bookimage{0.8}{info/images/09_output_window_avrdude_command.png}{Arduino IDE:s utskrift efter
uppladdningen, med avrdude-kommandot.}{studio:fig:avrdude}

\subsection{Gör om kommandot till ett återanvändbart verktyg}
\label{studio:5-3}
Kommandot ska delas upp i två delar: själva programmet (\textbf{Command}) och allt det får med sig
(\textbf{Arguments}). Gör det i Anteckningar (Notepad):
\begin{enumerate}
  \item Klistra in raden och ta bort alla citattecken.
  \item Sätt en radbrytning direkt efter \code{bin/avrdude}, och lägg till \code{.exe}, så att
    första raden slutar på \code{bin/avrdude.exe}.
\end{enumerate}

\bookimage{0.8}{info/images/10_notepad_path_pasted.png}{Anteckningar med avrdude-kommandot
inklistrat utan citattecken.}{studio:fig:pasted}

\bookimage{0.8}{info/images/11_notepad_exe_appended.png}{Anteckningar där första raden nu slutar
på avrdude.exe.}{studio:fig:exe}

\begin{enumerate}[start=3]
  \item Sätt en radbrytning till före \code{-Uflash}, så att resten hamnar på en tredje rad.
  \item Sätt tillbaka citattecken runt de två filsökvägarna: sökvägen till \code{avrdude.conf} på
    andra raden, och sökvägen till hexfilen på tredje.
\end{enumerate}

\bookimage{0.8}{info/images/13_notepad_quotes_added.png}{Kommandot på tre rader, med citattecken
runt de två filsökvägarna och Arduinos hexfil kvar på tredje raden.}{studio:fig:quotes}

\begin{enumerate}[start=5]
  \item Ersätt sökvägen till Arduinos hexfil, allt mellan citattecknen på tredje raden, med
    \code{$(ProjectDir)Debug\\$(TargetName).hex}. Microchip Studio byter ut de två variablerna mot
    det aktuella projektets katalog och namn, så samma kommando fungerar för alla dina projekt.
  \item Slå ihop de två sista raderna till en, så att det bara finns två rader kvar. Spara filen.
\end{enumerate}

\bookimage{0.8}{info/images/14_notepad_state_before_hex_replace.png}{Efter bytet i steg~5, med
sökvägen till projektets hexfil i Debug-katalogen på tredje raden.}{studio:fig:hexpath}

\bookimage{0.8}{info/images/12_notepad_arguments_split.png}{Samma tre rader utan den tomma raden,
och med citattecken också runt -PCOM4, vilket fungerar lika bra.}{studio:fig:variant}

\bookimage{0.8}{info/images/15_notepad_hex_path_replaced.png}{Efter steg~6, med två rader kvar och
filen sparad.}{studio:fig:saved}

Första raden är nu kommandot, andra raden argumenten.

\subsection{Lägg till verktyget i Microchip Studio}
\label{studio:5-4}
\begin{enumerate}
  \item Externa verktyg kräver den avancerade profilen: \textbf{Tools $\rightarrow$ Select
    Profile}, välj \textbf{Advanced} och klicka \textbf{Apply}.
\end{enumerate}

\bookimage{0.8}{info/images/19_select_profile_menu.png}{Microchip Studios meny Tools med Select
Profile markerat.}{studio:fig:profilemenu}

\bookimage{0.8}{info/images/20_advanced_profile_dialog.png}{Dialogen Select Profile med profilen
Advanced vald.}{studio:fig:advanced}

\begin{enumerate}[start=2]
  \item Välj \textbf{Tools $\rightarrow$ External Tools} och klicka \textbf{Add}.
  \item Döp verktyget till \code{Arduino} följt av COM-porten, till exempel \code{Arduino - COM4}.
  \item Klistra in första raden från Anteckningar i \textbf{Command} och andra raden i
    \textbf{Arguments}.
  \item Kryssa i \textbf{Use Output window} och klicka \textbf{OK}.
\end{enumerate}

\bookimage{0.8}{info/images/21_external_tools_menu.png}{Microchip Studios meny Tools med External
Tools markerat.}{studio:fig:externalmenu}

\bookimage{0.8}{info/images/22_external_tools_arduino_com4.png}{Dialogen External Tools, ifylld
med titel, kommando och argument för Arduino på COM4.}{studio:fig:com4}

Skärmbilderna är tagna i ett C-projekt. För ett assemblerprojekt är varje steg detsamma.

\subsection{Programmera och kontrollera}
\label{studio:5-5}
\begin{enumerate}
  \item Bygg projektet med \textbf{F7}. Hexfilen skapas i projektets katalog \code{Debug};
    kontrollera i Output-fönstret var den hamnade om programmeringen inte hittar den.
  \item Välj \textbf{Tools $\rightarrow$ Arduino - COM4}.
\end{enumerate}

\bookimage{0.8}{info/images/23_tools_menu_arduino_com4_selected.png}{Microchip Studios meny Tools
med det nya verktyget Arduino - COM4.}{studio:fig:toolcom4}

Om allt fungerar slutar utskriften med antalet byte som skrevs:

\begin{console}
Writing | ################################################## | 100% 0.05s

avrdude.exe: 194 bytes of flash written
...
avrdude.exe done. Thank you.
\end{console}

\bookimage{0.8}{info/images/24_output_window_flash_success.png}{Microchip Studios utskrift efter
en lyckad programmering av kortet.}{studio:fig:success}

\subsection{Om COM-porten ändras}
\label{studio:5-6}
Om kortet ansluts till en annan USB-port kan det få en annan COM-port, och programmeringen
misslyckas med ett fel i stil med:

\begin{textdrawing}
avrdude.exe: ser_open(): can't open device "\\.\COM4": Det går inte att hitta filen.
\end{textdrawing}

Kontrollera kortets aktuella port i Arduino IDE (\textbf{Tools $\rightarrow$ Port}). Lägg sedan
till ett andra verktyg i \textbf{Tools $\rightarrow$ External Tools}, likadant som det första men
med den nya porten i \textbf{Arguments}, till exempel \code{-PCOM3} i stället för \code{-PCOM4}.

\bookimage{0.8}{info/images/25_external_tools_menu_for_com3.png}{Microchip Studios meny Tools med
verktyget Arduino - COM4 och External Tools markerat.}{studio:fig:externalmenu2}

\bookimage{0.8}{info/images/26_external_tools_arduino_com3.png}{Dialogen External Tools, ifylld
för Arduino på COM3.}{studio:fig:com3}

\bookimage{0.8}{info/images/27_tools_menu_arduino_com3_selected.png}{Microchip Studios meny Tools
med verktygen för både COM3 och COM4.}{studio:fig:toolcom3}

\section{Vanliga fel}
\label{studio:6}

\begin{booktable}{@{}LL@{}}
  \thead{Symptom} & \thead{Trolig orsak} \\
  \midrule
  \textbf{Start Debugging} frågar efter ett verktyg & Simulatorn är inte vald, se
    \secref{studio:4-1}. \\
  Den gula pilen står kvar på samma rad & Programmet väntar i en loop, till exempel på en ingång
    som aldrig ändras. \\
  Ett register ändras inte som väntat & Den gula pilen visar nästa instruktion, inte den som just
    kördes. \\
  Stoppuret visar orimliga tider & \textbf{Frequency} står inte på 16~MHz. \\
  En port ändras inte i I/O-fönstret & Fel register (\code{PINB} i stället för \code{PORTB}),
    eller \code{DDRB} är inte satt. \\
  \code{avrdude} hittar inte hexfilen & Projektet är inte byggt, eller hexfilen ligger i en annan
    katalog än \code{Debug}. \\
  \code{ser_open(): can't open device} & Fel COM-port, eller kortet är inte anslutet, se
    \secref{studio:5-6}. \\
\end{booktable}
